% 최종 축약본. chapter4.md와 같은 본문에서 생성.
\subsection{오류 유형 및 관측 사례}
\label{sec:ch4_errors}

Jev는 상위 후보에서 정답 근거가 탈락하는 경우를 줄였지만, 이후 선택되지 않는 근거도 남았다. 질문마다 필요한 고유 정답 근거를 따로 세면 총 443건이다. 검색 자료에서 동일 문자열을 찾을 수 없는 근거 33건과, 자료에는 있으나 검색 후보에 없는 75건은 두 조건에서 같았다. 검색 후보에는 있으나 재정렬 후 상위 10개에서 탈락한 근거는 CE 42건, Jev 13건이었다. 상위 10개에 남았으나 최종 선택되지 않은 근거는 60건과 72건이었고, 최종 선택은 233건과 250건이었다. 이 집계는 같은 근거가 여러 질문에 필요하면 각각 세므로, 전체 고유 근거 386개와 다르다.

개선·악화·지표 불일치 집단에서 근거 F1 차이의 크기가 각 집단의 중앙값에 가장 가까운 사례를 하나씩 골라 선택 근거와 답변의 변화를 살펴보았다(표~\ref{tab:ch4_cases}).

\begin{table}[htbp]
\centering
\caption{관측 사례의 점수 변화 (CE $\to$ Jev)}
\label{tab:ch4_cases}
\small
\setlength{\tabcolsep}{3pt}
\begin{tabular}{@{}lrrr@{}}
\hline
사례 & 근거 F1 & Token F1 & LF-Eval \\
\hline
A & 50.00 $\to$ 100.00 & 69.09 $\to$ 86.79 & 65.29 $\to$ 99.18 \\
B & 80.00 $\to$ 40.00 & 32.79 $\to$ 29.09 & 71.79 $\to$ 52.67 \\
C & 40.00 $\to$ 0.00 & 18.60 $\to$ 25.88 & 31.10 $\to$ 30.55 \\
\hline
\end{tabular}
\end{table}

사례 A는 식품 생산실적과 회수계획의 보고 상대방을 묻는 질문이다. CE는 식품위생법 제42조제2항과 생산실적 관련 시행규칙을 선택하였다. Jev는 제42조제2항을 유지하면서 회수계획 관련 정답 근거인 제45조제1항도 확보하였다. 답변에는 시장·군수·구청장이 포함되었고 두 답변 점수도 높아졌다.

사례 B는 여러 당사자에 대한 복수 청구 절차를 묻는 질문이다. CE는 정답 근거인 민사소송법 제79조제1항과 제253조를 모두 선택하였다. Jev에서는 제79조제1항이 상위 후보에서 탈락하고 제83조제1항이 선택되었다. 답변도 독립당사자참가에서 공동소송참가 관련 내용으로 달라졌으며, 근거와 답변 점수는 모두 낮아졌다.

사례 C는 사면과 복수 범죄의 처리를 묻는 질문이다. 정답 근거 3개 중 하나는 검색 자료에서 동일 문자열을 찾을 수 없었고, 다른 하나는 검색 후보에 없었다. 후보 목록도 2개여서 목록마다 하나를 선택하는 방식으로는 정답 근거 3개를 모두 확보할 수 없었다. CE는 정답 근거 하나를 골랐지만 Jev는 고르지 못했다. 그럼에도 Token F1은 높아져, 정답과 표현이 더 비슷한 답변이 반드시 더 나은 근거 선택을 뜻하지 않음을 보여준다.

이 사례들은 자동 관측 결과이며 법률 전문가의 판정은 완료되지 않았다. 특히 상위 후보에 있는 근거를 선택하지 않은 경우에는 목록당 하나만 고르는 제약도 작용할 수 있다. 정답과 문자열이 다른 근거의 법률적 타당성 역시 별도 검토가 필요하다. 상세 평가와 집계·사례 선정 절차는 보충자료에 제시하였다.
